\section{Introduction}\label{sec:introduction}
The reconstruction problem asks whether many distant, noisy descendants
retain information about their common ancestor. On a tree, the growth in
the number of descendants competes with the loss of information across
each edge. For a symmetric channel with nontrivial eigenvalue $\lambda$
and mean offspring number $d$, the Kesten--Stigum condition
$d\lambda^2>1$ gives reconstruction. Whether this condition is also
necessary depends on the channel. We determine the answer for three
states at every branching parameter in two basic tree families.
Through existing tree-to-graph and algorithmic results, the Poisson
theorem also determines the weak-recovery threshold for the symmetric
three-community sparse stochastic block model with positive connectivity
parameters and mean degree greater than one.

\subsection{Model and result}
Let the root spin $\sigma_\rho$ be uniform on $\{1,2,3\}$. Along each
edge, independently conditional on its parent spin, transmit through
\begin{equation}\label{eq:channel}
 K_\lambda(i,j)=
 \begin{cases}
 (1+2\lambda)/3,&i=j,\\
 (1-\lambda)/3,&i\ne j,
 \end{cases}
 \qquad -\tfrac12\le\lambda\le1.
\end{equation}
The tree is either the rooted $b$-ary tree, with an integer $b\ge2$
children at every vertex, or a Galton--Watson tree with independent
$\operatorname{Pois}(d)$ offspring counts, where $d>1$.
For the regular model set $d=b$. The tree is independent of the root
spin. The observation at depth $\ell$ consists of the unlabelled tree
$T_{\le\ell}$ through that depth together with the spins attached to its
vertices $\partial T_\ell$ at depth $\ell$; neither interior spins nor
an ordering of siblings are observed. Define
\begin{equation}\label{eq:advantage}
 a_\ell=\E\left\|
 \PP(\sigma_\rho=\cdot\mid T_{\le\ell},\sigma_{\partial T_\ell})
 -\operatorname{Unif}\{1,2,3\}\right\|_{\mathrm{TV}}.
\end{equation}
In the Poisson case this expectation includes both the tree and the
spins, \emph{without conditioning on survival}. Reconstruction means
$\lim_{\ell\to\infty}a_\ell>0$; non-reconstruction means that this
limit is zero. The limit exists by information degradation.

\begin{theorem}[Exact threshold]\label{thm:threshold}
For every integer $b\ge2$ in the regular model and every real $d>1$
in the Poisson model, reconstruction occurs if and only if
\[
 d\lambda^2>1,
\]
with $d=b$ in the regular case. At every admissible equality
$d\lambda^2=1$, there is non-reconstruction.
\end{theorem}
Thus the positive threshold is $\lambda=d^{-1/2}$. On the negative
side reconstruction occurs precisely when $d>4$ and
$-1/2\le\lambda<-d^{-1/2}$. In particular, the three-coloring channel
$\lambda=-1/2$ is non-reconstructible at branching parameter four in
both models.

\subsection{History and significance}
The branching-process work of Kesten and Stigum~\cite{KestenStigum1966}
underlies the classical second-moment reconstruction bound. Broadcasting
and its relation to the Ising model were developed by Evans, Kenyon,
Peres, and Schulman~\cite{EvansEtAl2000}. For more than two states,
nonlinear information in the boundary can change the threshold. This
possibility occurs for symmetric channels with five or more states
\cite[Theorem~1.2]{Sly2011}; the three-state assertion therefore needs
more than the second-moment bound.

M\'ezard and Montanari~\cite[Conjecture~2 and the following paragraph]{MezardMontanari2006}
predicted sharpness of the Kesten--Stigum bound for three-state symmetric
channels of either sign up to a branching cutoff, and proposed the
extension to all branching parameters. See also Conjecture~1 of
Sly~\cite{Sly2011} for the all-degree regular-tree formulation.
Sly proved this for all sufficiently large regular branching parameters,
including non-reconstruction at equality
\cite[Theorem~1.1]{Sly2011}. Mossel, Sly, and Sohn established the
corresponding large-mean Poisson-tree result, for both signs and including
equality \cite[Theorem~1.12]{MosselSlySohn2025}.
At small degrees, Bhatnagar and Maneva developed finite representations
of posterior laws and proved, among other bounds, non-reconstruction
for three-colorings on the regular ternary tree
\cite[Theorem~12]{BhatnagarManeva2011}.
Theorem~\ref{thm:threshold} removes the large-branching restriction in
these two tree families. It resolves the three-state regular-tree
prediction in its all-degree form and includes the low-branching and
coloring endpoints for the observed Poisson tree.

For the symmetric stochastic block model, the cavity-method analysis of
Decelle, Krzakala, Moore, and Zdeborov\'a predicted the Kesten--Stigum
transition in the assortative three-community case
\cite[Section~IV.A]{DecelleEtAl2011}.
Ricci-Tersenghi, Semerjian, and Zdeborov\'a gave further evidence for the
three-state prediction for either sign, through local analysis of the
posterior recursion and numerical cavity solutions
\cite[Sections~III.D.1 and~IV.H]{RicciTersenghiEtAl2019}.
Mossel, Sly, and Sohn proved information-theoretic impossibility at and
below the Kesten--Stigum bound for three or four communities when the
mean degree is sufficiently large
\cite[Theorem~1.5]{MosselSlySohn2025}.
Their tree-to-graph transfer has no large-degree restriction
\cite[Theorem~2.1]{MosselSlySohn2025}.
Together with the above-threshold algorithm of Abbe and Sandon
\cite[Corollary~1 and Lemma~1]{AbbeSandon2018}, the Poisson part of our
theorem therefore gives the exact three-community weak-recovery threshold
stated in Corollary~\ref{cor:sbm} below.

A related three-state model with branch-dependent couplings on a binary
Cayley tree is studied by Mukhamedov and Ak{\i}n~\cite{MukhamedovAkin2026}.
Their work concerns disordered translation-invariant Gibbs measures and
their extremality; our theorem concerns the same homogeneous symmetric
channel on every edge of the regular or observed Poisson tree.

\subsection{Proof strategy and its precedents}
For a root posterior $p=(p_1,p_2,p_3)$, write $m_i=3p_i$ and identify
the three spin labels with $1,\omega,\omega^2$, where
$\omega=e^{2\pi\mathrm{i}/3}$. The complex posterior mean and its two
basic invariants are
\[
 \zeta=p_1+\omega p_2+\omega^2p_3,\qquad
 x=|\zeta|^2,\qquad y=\Re(\zeta^3).
\]
The posterior recursion passes each child experiment through one edge
and combines the resulting conditionally independent branch observations.
Section~\ref{fnd:section} gives this recursion and proves that the
limiting posterior law is a fixed point, including at the coloring
endpoint. Non-reconstruction is equivalent to $\E x=0$ under that law.
A linear estimator proves reconstruction above the threshold. Channel
and offspring degradation reduce the converse to critical equality and
the coloring case at branching parameter four.

At positive equality, Section~\ref{sec:positive} compares two expressions
for the change of an entropy functional across the recursion.
The expected symmetrized-entropy identity goes back to Formentin and
K\"ulske~\cite[Proposition~3.3 and equation~(15)]{FormentinKulske2009Entropy}.
Our functional combines entropy and symmetrized entropy with a quadratic
centered-logarithm term. The fixed-point identities express its change
through the correlation generated by
combining branches, while a radial inequality bounds it below by a
positive multiple of $\E x^2$. Moment estimates make these two bounds
incompatible at a nonuniform limiting fixed point. Appendix~\ref{rad:section}
proves the radial inequality analytically. Related entropy methods
include the sharp nonlinear information-contraction inequalities for the
Potts channel and non-reconstruction criteria of Gu and
Polyanskiy~\cite{GuPolyanskiy2023}.

The negative argument separates small and large posterior moments.
Sly's large-branching proof already combines a moment estimate near the
uniform posterior with a global Gaussian contraction, including a
rigorous numerical verification
\cite[Section~1.5, Lemmas~4.7--4.8, and Appendix~A]{Sly2011}.
Here Section~\ref{sec:local} proves a uniform quantitative gap: at every
negative equality, any nonuniform symmetric fixed point satisfies
$\E x>1/20$. Pair inequalities for $x$, $y$, and $x^2$ yield this gap
through a moment bootstrap. Section~\ref{sec:weighted} then uses weighted
logarithm identities to exclude the remaining fixed points away from
the coloring channel.

Near coloring, Section~\ref{grid:section} constructs finite experiments
that are more informative than the true boundary observation. It
suffices for one such upper experiment to have $\E x<1/20$: the true
limiting fixed point then falls below the gap and must be uniform.
The upper recursion need not itself converge to the uniform experiment.
Bhatnagar and Maneva introduced finite surveys of posterior laws and
combined certified finite recursion with local contraction
\cite[Sections~3.1--3.3 and Theorem~12]{BhatnagarManeva2011}.
Their construction includes completion by perfectly informative messages
for omitted offspring tails and downward-rounded masses
\cite[Section~3.3, implementation remarks~1 and~4]{BhatnagarManeva2011}.
We use this finite-experiment approach with the gap above and explicit
three-state refinement kernels. Gu's recent binary-hypertree work proves
a related rigorous population-dynamics procedure and proposes a Potts
extension based on degradation and local non-reconstruction
\cite[Sections~4 and~7]{Gu2026Hypertrees}.

The quantitative radial bound, uniform fixed-point gap, and exact
certificates supply the estimates needed throughout the parameter range.
The negative argument is computer-assisted: the polynomial, interval,
and finite-experiment programs are printed in
Appendix~\ref{app:certificates}. Their mathematical interpretation and
arithmetic bounds are proved in the text. All decimal constants denote
exact terminating rationals, and no simulation estimate enters a
certification decision. Section~\ref{sec:completion} assembles the tree
classification; Table~\ref{tab:coverage} there records the parameter
ranges of its analytic and finite arguments.

\subsection{A consequence for community detection}\label{sec:sbm}
We now use $a$ and $b$ for within- and between-community connectivity
parameters, independently of the notation for regular branching above.
Fix $a,b>0$ with $d=(a+2b)/3>1$. For all sufficiently large $n$, choose
$\sigma_1,\ldots,\sigma_n$ independently and uniformly from $\{1,2,3\}$.
Conditional on these labels, form a graph $G_n$ on $\{1,\ldots,n\}$ by
including each unordered pair independently, with probability $a/n$ if
its endpoints have the same label and $b/n$ otherwise. For two labelings
$\sigma,\widehat\sigma\in\{1,2,3\}^n$, define their overlap by
\[
 \operatorname{ov}(\sigma,\widehat\sigma)
 =\max_{\pi\in S_3}\frac1n\sum_{v=1}^n
   \mathbf1_{\{\widehat\sigma_v=\pi(\sigma_v)\}}.
\]
The permutation accounts for the fact that the names of the communities
cannot be recovered from the graph.

\begin{corollary}[Exact weak-recovery threshold for three communities]
\label{cor:sbm}
In the model above, the following statements hold.
\begin{enumerate}[label=(\roman*)]
\item If $(a-b)^2\le 3(a+2b)$, then every sequence of estimators
$\widehat\sigma_n$ using only $G_n$ and independent randomization satisfies
\[
 \lim_{n\to\infty}\E\operatorname{ov}(\sigma,\widehat\sigma_n)=\frac13.
\]
\item If $(a-b)^2>3(a+2b)$, there are $\varepsilon>0$ and an algorithm,
allowed to depend on the fixed parameters $a,b$, which runs in
$O(n\log n)$ time and satisfies
\[
 \PP\!\left(\operatorname{ov}(\sigma,\widehat\sigma_n)
              \ge\frac13+\varepsilon\right)\longrightarrow1.
\]
\end{enumerate}
Thus weak recovery is possible exactly under the strict inequality,
and is information-theoretically impossible at equality.
\end{corollary}

The proof of Corollary~\ref{cor:sbm}, including the passage from averaged
tree non-reconstruction to the fixed-tree hypothesis of the transfer
theorem, is given in Section~\ref{sec:completion}.
